%=============================================================================
% ALL ALPHA EQUATIONS IN DFD
% Comprehensive cross-reference of every equation involving the fine structure
% constant, stiffness ratios, Wilson ratios, and related quantities
% Compiled: 2026-03-22
%=============================================================================

\documentclass[11pt]{article}
\usepackage{amsmath,amssymb,amsthm}
\usepackage{geometry}
\geometry{margin=1in}
\usepackage{tcolorbox}
\usepackage{booktabs}

\newcommand{\alphafine}{\alpha}

\title{Complete Catalogue of $\alpha$-Related Equations in DFD}
\author{Cross-Reference Compilation}
\date{22 March 2026}

\begin{document}
\maketitle

\tableofcontents
\newpage

%=============================================================================
\section{The Definition of $\alpha$}
%=============================================================================

From the Ab Initio paper (Eq.~1) and section\_quantum.tex:

\begin{equation}\tag{AbInit-1}
\alpha \equiv \frac{e^2}{4\pi\epsilon_0\hbar c} \approx 0.0072973525693\ldots \approx \frac{1}{137.036}
\end{equation}

%=============================================================================
\section{The Complete Derivation Chain: From Topology to $\alpha$}
%=============================================================================

The derivation chain (Ab Initio paper, Eq.~A6; section\_quantum.tex):
\begin{equation}\tag{Chain}
\text{SM hypercharge} \;\to\; q_1 = 3 \;\to\; a = 9 \;\to\; k_{\max} = 60 \;\to\; \alpha = 1/137.
\end{equation}

%-----------------------------------------------------------------------------
\subsection{Constraint 1: The UV Cutoff $k_{\max} = 60$ (Bridge Lemma)}
%-----------------------------------------------------------------------------

\paragraph{Bridge Lemma (Ab Initio Eq.~14, Eq.~A1, Eq.~A5; section\_quantum.tex):}
\begin{equation}\tag{BL-1}
k_{\max} = \chi(\mathbb{C}P^2, \mathcal{O}(9) \oplus \mathcal{O}^{\oplus 5}) = \binom{11}{2} + 5 = 55 + 5 = 60.
\end{equation}

Formal index form (Ab Initio Eq.~A1):
\begin{equation}\tag{BL-2}
k_{\max} := \mathrm{Index}(D_{\mathbb{C}P^2} \otimes E) = \chi(\mathbb{C}P^2, E) = 60.
\end{equation}

Hirzebruch--Riemann--Roch components (Ab Initio Eqs.~A2--A5):
\begin{align}
\mathrm{Index}(D_{\mathbb{C}P^2} \otimes E) &= \chi(\mathbb{C}P^2, E), \tag{BL-3} \\
\chi(\mathcal{O}(9)) &= \binom{11}{2} = 55, \tag{BL-4} \\
\chi(\mathcal{O}) &= 1, \tag{BL-5} \\
k_{\max} = \chi(E) &= \chi(\mathcal{O}(9)) + 5\chi(\mathcal{O}) = 55 + 5 = 60. \tag{BL-6}
\end{align}

\paragraph{Uniqueness constraint (Ab Initio, App.~A.3):}
With $E = \mathcal{O}(a) \oplus \mathcal{O}^{\oplus n}$ and $\chi(E) = 60$, the unique
minimal-padding solution is $(a,n) = (9,5)$:
\begin{equation}\tag{BL-7}
\binom{a+2}{2} + n = 60, \quad a \leq 9 \;\;\text{(since } \binom{12}{2} = 66 > 60\text{)}.
\end{equation}

\paragraph{Icosahedral confirmation:}
\begin{equation}\tag{BL-8}
k_{\max} = |A_5| = 60 \quad \text{(McKay correspondence: } 2I \subset SU(2) \leftrightarrow E_8\text{)}.
\end{equation}

\paragraph{$E_8$ echo:}
\begin{equation}\tag{BL-9}
k_{\max} = \mathrm{roots}(E_8)/4 = 240/4 = 60.
\end{equation}

%-----------------------------------------------------------------------------
\subsection{Constraint 2: Microsector Vacuum Sets $\beta_{U(1)}$}
%-----------------------------------------------------------------------------

\paragraph{Chern--Simons partition function on $S^3$ (Ab Initio Eq.~3):}
\begin{equation}\tag{CS-1}
Z_{SU(2)_k}(S^3) = \sqrt{\frac{2}{k+2}}\sin\!\left(\frac{\pi}{k+2}\right).
\end{equation}

\paragraph{Shifted level (Ab Initio Eq.~4):}
\begin{equation}\tag{CS-2}
k_{\mathrm{eff}} \equiv k + 2.
\end{equation}

\paragraph{Microsector weight (Ab Initio Eq.~5):}
\begin{equation}\tag{CS-3}
w(k) = |Z_{SU(2)_k}(S^3)|^2 = \frac{2}{k+2}\sin^2\!\left(\frac{\pi}{k+2}\right).
\end{equation}

\paragraph{Vacuum expectation (Ab Initio Eq.~15):}
\begin{equation}\tag{CS-4}
\boxed{\langle k_{\mathrm{eff}}\rangle = \frac{\sum_{k=0}^{59}(k+2)\,w(k)}{\sum_{k=0}^{59} w(k)} = 3.7969 \approx 3.80.}
\end{equation}

\paragraph{Dictionary entry (Ab Initio Eq.~16):}
\begin{equation}\tag{CS-5}
\boxed{\beta_{U(1)} = \langle k_{\mathrm{eff}}\rangle_{k_{\max}=60} = 3.80.}
\end{equation}

\paragraph{Sensitivity to $k_{\max}$ (Ab Initio Table I):}
\begin{center}
\begin{tabular}{ccc}
\toprule
$k_{\max}$ & $\langle k+2\rangle$ & $\alpha$ result \\
\midrule
50 & 3.77 & $1/137$ (+1.3\%) \\
\textbf{60} & \textbf{3.80} & $\mathbf{1/137}$ \textbf{(+0.5\%)} \\
$\infty$ & 3.94 & $1/303$ ($-55\%$, \textbf{ruled out}) \\
\bottomrule
\end{tabular}
\end{center}

%-----------------------------------------------------------------------------
\subsection{Constraint 3: The Lattice Ratio from Topology}
%-----------------------------------------------------------------------------

\paragraph{Lattice ratio (Ab Initio Eq.~17):}
\begin{equation}\tag{LR-1}
\boxed{\frac{\beta_{SU(2)}}{\beta_{U(1)}} = \frac{n_2}{n_1} \times N_{\mathrm{gen}} = 2 \times 3 = 6.}
\end{equation}

This combines:
\begin{itemize}
\item Frame Stiffness Theorem: $n_2/n_1 = \kappa_{SU(2)}/\kappa_{U(1)} = 2$
\item Index theorem on $\mathbb{C}P^2$: $N_{\mathrm{gen}} = 3$
\end{itemize}

\paragraph{Resulting $\beta_{SU(2)}$ (Ab Initio Eq.~18):}
\begin{equation}\tag{LR-2}
\beta_{SU(2)} = 6 \times 3.80 = 22.80.
\end{equation}

%-----------------------------------------------------------------------------
\subsection{Constraint 4: DFD Stiffness Ratio (Frame Stiffness Theorem)}
%-----------------------------------------------------------------------------

\paragraph{Stiffness proportionality (Ab Initio Eq.~8):}
\begin{equation}\tag{FST-1}
\kappa_r = n_r \kappa_0,
\end{equation}
where $n_{U(1)} = 1$, $n_{SU(2)} = 2$, $n_{SU(3)} = 3$ are complex dimensions of internal mode subspaces.

\paragraph{Stiffness ratio (Ab Initio Eq.~9):}
\begin{equation}\tag{FST-2}
\boxed{\frac{\kappa_{U(1)}}{\kappa_{SU(2)}} = \frac{n_1}{n_2} = \frac{1}{2}.}
\end{equation}

Lattice measurement: $\kappa_{U(1)}/\kappa_{SU(2)} = 0.495 \pm 0.020$.

%=============================================================================
\section{Gauge Coupling Extraction}
%=============================================================================

\paragraph{Stiffness functional (Ab Initio Eq.~7; section\_quantum.tex):}
\begin{equation}\tag{GC-1}
\mathcal{L}^{(r)}_{\mathrm{stiff}} = -\frac{\kappa_r}{2}\,\mathrm{Tr}\,F^{(r)}_{ij}F^{(r)ij}.
\end{equation}

Canonical normalization: $g_r \propto \kappa_r^{-1/2}$.

\paragraph{Berry connection (Ab Initio Eq.~6):}
\begin{equation}\tag{GC-2}
A_i^{(r)} = i\,\Xi_r^\dagger \partial_i \Xi_r.
\end{equation}

\paragraph{Wilson normalization (Ab Initio Eq.~12):}
\begin{equation}\tag{GC-3}
g_1^2 = \frac{1}{\kappa_{U(1)}}, \qquad g_2^2 = \frac{4}{\kappa_{SU(2)}}.
\end{equation}

\paragraph{Electroweak mixing (Ab Initio Eq.~10):}
\begin{equation}\tag{GC-4}
\frac{1}{e^2} = \frac{1}{g_1^2} + \frac{1}{g_2^2}.
\end{equation}

\paragraph{Wilson-normalized $\alpha_W$ (Ab Initio Eq.~13):}
\begin{equation}\tag{GC-5}
\boxed{\alpha_W = \frac{e^2}{4\pi} = \frac{(1/\kappa_{U(1)})(4/\kappa_{SU(2)})}{(1/\kappa_{U(1)}) + (4/\kappa_{SU(2)})} \cdot \frac{1}{4\pi}.}
\end{equation}

\paragraph{Lattice input--output distinction (Ab Initio Eq.~11):}
\begin{equation}\tag{GC-6}
\beta_{U(1)} = \frac{1}{g_1^2} \approx 3.80 \;\;\text{(input)}, \qquad \kappa_{U(1)} \approx 7.25 \;\;\text{(measured)}.
\end{equation}

%=============================================================================
\section{Lattice Action}
%=============================================================================

\paragraph{DFD micro-action (Ab Initio Eq.~19):}
\begin{equation}\tag{LA-1}
S = \sum_x [-\log w(k_x)] + \frac{K_\psi}{2}\sum_{\langle xy\rangle}(\psi_x - \psi_y)^2 - \sum_p \beta\,e^{-\psi_p}\cos(\theta_p + \theta_{\mathrm{bg}}\,\Omega_p).
\end{equation}

\paragraph{Kappa extraction (Ab Initio Eq.~20):}
\begin{equation}\tag{LA-2}
\kappa = \frac{F''(0)}{V}, \qquad F''(0) = \langle S''\rangle - \langle (S')^2\rangle + \langle S'\rangle^2.
\end{equation}

%=============================================================================
\section{Convention-Locked $\alpha$ (Full DFD v10.8, section\_alpha\_locked.tex)}
%=============================================================================

%-----------------------------------------------------------------------------
\subsection{Operator and Truncation}
%-----------------------------------------------------------------------------

\paragraph{Connection Laplacian on $\mathbb{C}P^2 \times S^3$:}
\begin{equation}\tag{CL-1}
P = -g^{ij}\nabla_i \nabla_j.
\end{equation}

\paragraph{Spectral action with plateau cutoff:}
\begin{equation}\tag{CL-2}
S = \mathrm{Tr}\,f(P/\Lambda^2), \qquad f_{-2} = f_{-4} = \cdots = 0.
\end{equation}

\paragraph{Toeplitz truncation dimension:}
\begin{equation}\tag{CL-3}
d = \dim H^0(\mathbb{C}P^1, \mathcal{O}(m)) = m + 1 = k + 4.
\end{equation}

For $k = 60$: $d = 64$.

\paragraph{Spin$^c$ origin of the +3 shift:}
\begin{equation}\tag{CL-4}
K_{\mathbb{C}P^2} = \mathcal{O}(-3) \quad\Rightarrow\quad L_{\det} = K^{-1} = \mathcal{O}(3).
\end{equation}

\paragraph{Spectral cutoff:}
\begin{equation}\tag{CL-5}
\Lambda^3 = k \cdot \frac{d-1}{d} = k \cdot \frac{k+3}{k+4}.
\end{equation}

For $k = 60$: $\Lambda^3 = 60 \times 63/64 = 885.9375$.

%-----------------------------------------------------------------------------
\subsection{The Forced Microsector Fork}
%-----------------------------------------------------------------------------

\paragraph{Branch A: Regular-module microsector (survives):}
\begin{align}
\mathcal{H}_F &:= A = M_d(\mathbb{C}), \tag{Fork-A1} \\
\mathrm{tr}_{\mathrm{dem}}(\cdot) &:= \frac{1}{d^2}\,\mathrm{Tr}_{\mathcal{H}_F}(\cdot), \tag{Fork-A2} \\
\varepsilon_{\mathrm{adj}}^{(A)} &= \frac{d^2}{d^2 - 1} = \frac{4096}{4095} = 1.000244\ldots \tag{Fork-A3}
\end{align}

\paragraph{Branch B: Fermion-rep microsector (falsified):}
\begin{equation}\tag{Fork-B1}
\varepsilon_{\mathrm{adj}}^{(B)} = \frac{d^2 - 1}{d^2} = \frac{4095}{4096} = 0.999756\ldots
\end{equation}

%-----------------------------------------------------------------------------
\subsection{Complete Derivation Chain Components}
%-----------------------------------------------------------------------------

\begin{center}
\begin{tabular}{llll}
\toprule
Component & Value & Source & Status \\
\midrule
$K_{\mathbb{C}P^2}$ & $\mathcal{O}(-3)$ & Algebraic geometry & Rigorous \\
$L_{\det} = K^{-1}$ & $\mathcal{O}(3)$ & Spin$^c$ structure & Rigorous \\
$d = k + 4$ & 64 & $\dim H^0(\mathcal{O}(k+3))$ & Rigorous \\
$(d-1)/d$ & 63/64 & Traceless projection & Derived \\
$N_{\mathrm{species}}$ & 7 & SM SU(2) components & SM content \\
$\mathrm{Tr}(Y^2)$ & 10 & SM hypercharges & SM content \\
$g_F$ & 8 & Spectral triple ($J \times \gamma \times \mathbb{C}$) & Derived \\
$w = N_{\mathrm{species}}/(g_F \cdot \mathrm{Tr}(Y^2))$ & 7/80 & Hypercharge weighting & Derived \\
$\varepsilon_{\mathrm{adj}}$ & 4096/4095 & Regular-module trace & Forced \\
\midrule
$\alpha^{-1}$ & 137.03599985 & All above combined & $< 0.01$ ppm \\
\bottomrule
\end{tabular}
\end{center}

%-----------------------------------------------------------------------------
\subsection{The Result}
%-----------------------------------------------------------------------------

\begin{tcolorbox}[colback=green!5!white, colframe=green!50!black, title=Convention-Locked Result]
\begin{equation}\tag{RESULT-1}
\boxed{\alpha^{-1} = 137.03599985 \quad \text{(residual: } -0.006\text{ ppm)}}
\end{equation}

Experimental: $\alpha^{-1}_{\mathrm{exp}} = 137.035999084$.

Branch B gives $\alpha^{-1} = 137.03014445$ (residual: $+42.7$ ppm) --- \textbf{falsified}.
\end{tcolorbox}

\paragraph{Lattice-level result (section\_quantum.tex):}
\begin{equation}\tag{RESULT-2}
\boxed{\alpha^{-1} = 137.036 \pm 0.5}
\end{equation}

%=============================================================================
\section{The $\alpha$-Relations (section\_alpha\_relations.tex)}
%=============================================================================

%-----------------------------------------------------------------------------
\subsection{Relation I: Self-Coupling $k_a$}
%-----------------------------------------------------------------------------

\begin{equation}\tag{AR-1}
\boxed{k_a = \frac{3}{8\alpha} \approx 51.4}
\end{equation}

Decomposition:
\begin{equation}\tag{AR-2}
k_a = N_{\mathrm{gen}} \times C_{\mathrm{loop}} \times \frac{1}{\alpha} = 3 \times \frac{1}{8} \times \frac{1}{\alpha}.
\end{equation}

%-----------------------------------------------------------------------------
\subsection{Relation II: EM Threshold}
%-----------------------------------------------------------------------------

\begin{equation}\tag{AR-3}
\boxed{\eta_c = \alpha \times \sin^2\theta_W \approx \frac{\alpha}{4}}
\end{equation}

%-----------------------------------------------------------------------------
\subsection{Relation III: Clock Coupling}
%-----------------------------------------------------------------------------

\begin{equation}\tag{AR-4}
\boxed{k_\alpha = \alpha \times a_e = \frac{\alpha^2}{2\pi}}
\end{equation}

where $a_e = \alpha/(2\pi)$ is the Schwinger anomalous magnetic moment.

Source: appendix\_P\_kalpha\_MR.tex, Theorem statement.

%-----------------------------------------------------------------------------
\subsection{Relation IV: MOND Scale (Derived)}
%-----------------------------------------------------------------------------

From variational stationarity (section\_alpha\_relations.tex; appendix\_N\_mu\_derivation.tex):
\begin{equation}\tag{AR-5}
k_a \times a_0^2 = \frac{3}{2}(cH_0)^2,
\end{equation}

therefore:
\begin{equation}\tag{AR-6}
a_0^2 = \frac{3(cH_0)^2}{2 \times \frac{3}{8\alpha}} = 4\alpha(cH_0)^2,
\end{equation}

giving:
\begin{equation}\tag{AR-7}
\boxed{a_0 = 2\sqrt{\alpha}\,cH_0.}
\end{equation}

Numerical: $a_0^{\mathrm{derived}} = 1.13 \times 10^{-10}$ m/s$^2$ vs.\ $a_0^{\mathrm{obs}} = (1.20 \pm 0.26) \times 10^{-10}$ m/s$^2$.

%-----------------------------------------------------------------------------
\subsection{Consistency Checks}
%-----------------------------------------------------------------------------

\begin{align}
\eta_c \times k_a &= \frac{\alpha}{4} \times \frac{3}{8\alpha} = \frac{3}{32} \quad\text{($\alpha$-independent)}, \tag{CC-1} \\
k_a \times a_0^2 &= \frac{3}{8\alpha} \times 4\alpha(cH_0)^2 = \frac{3}{2}(cH_0)^2 \quad\text{($\alpha$-independent)}, \tag{CC-2} \\
\frac{k_\alpha}{\alpha \times a_e} &= \frac{\alpha^2/(2\pi)}{\alpha \times \alpha/(2\pi)} = 1 \quad\text{(exact)}. \tag{CC-3}
\end{align}

%=============================================================================
\section{Electroweak Hierarchy: $v$ from $\alpha$ (section\_quantum.tex)}
%=============================================================================

\begin{equation}\tag{EH-1}
\boxed{v = M_P \times \alpha^8 \times \sqrt{2\pi}}
\end{equation}

Numerical verification:
\begin{align}
M_P &= 1.221 \times 10^{19}\;\text{GeV}, \tag{EH-2} \\
\alpha^8 &= (1/137.036)^8 = 8.04 \times 10^{-18}, \tag{EH-3} \\
\sqrt{2\pi} &= 2.507, \tag{EH-4} \\
v_{\mathrm{pred}} &= 246.09\;\text{GeV} \quad\text{(exp: 246.22 GeV)}. \tag{EH-5}
\end{align}

%=============================================================================
\section{Fermion Mass Formula (section\_quantum.tex)}
%=============================================================================

\begin{equation}\tag{FM-1}
\boxed{m_f = A_f \cdot \alpha^{n_f} \cdot \frac{v}{\sqrt{2}}}
\end{equation}

where $\alpha = 1/137.036$, $v/\sqrt{2} = 174.1$ GeV, $n_f$ is the sector-dependent
exponent from $\mathbb{C}P^2$ coupling path, and $A_f$ is a rational prefactor.

%=============================================================================
\section{Generation Count (section\_quantum.tex)}
%=============================================================================

\begin{equation}\tag{GEN-1}
N_{\mathrm{gen}} = |k_3 \cdot k_2 \cdot q_1|.
\end{equation}

With $k_3 = k_2 = q_1 = 1$: $N_{\mathrm{gen}} = 1$ per flux unit.
With minimal hypercharge twist $q_1 = 3$: contributes to generation count.

%=============================================================================
\section{Lattice Monte Carlo Results at $(\beta_{U(1)}, \beta_{SU(2)}) = (3.80, 22.80)$}
%=============================================================================

\paragraph{Ab Initio paper headline results:}

\begin{center}
\begin{tabular}{cccc}
\toprule
$L$ & $n$ & $\alpha_W$ (mean) & $\Delta\alpha/\alpha$ \\
\midrule
6 & 5 & \textbf{0.007297} & $-0.00\%$ \\
8 & 5 & 0.007322 & $+0.34\%$ \\
10 & 4 & 0.007361 & $+0.88\%$ \\
12 & 2 & \textbf{0.007291} & $-0.08\%$ \\
16 & 9$^\dagger$ & 0.007380 & $+1.13\%$ \\
\bottomrule
\end{tabular}
\end{center}

Physical value: $\alpha_{\mathrm{phys}} = e^2/(4\pi) = 1/(4\pi \times 137.036) = 0.007297$.

%=============================================================================
\section{Wilson Ratio Verification (Ab Initio Table III, Fig.~2)}
%=============================================================================

Only the DFD-predicted ratio of 6 yields $\alpha = 1/137$. Ten ratios tested (3--9
including fractional 5.5, 6.25, 6.5); all except exactly 6 fail.

\begin{center}
\begin{tabular}{cccc}
\toprule
$\beta_{SU(2)}/\beta_{U(1)}$ & $\beta_{SU(2)}$ & $\alpha_W$ & Deviation \\
\midrule
3 & 11.40 & 0.008907 & $+22.1\%$ \\
4 & 15.20 & 0.008234 & $+12.8\%$ \\
5 & 18.85 & 0.008005 & $+9.7\%$ \\
5.5 & 20.90 & 0.007549 & $+3.5\%$ \\
\textbf{6} & \textbf{22.80} & \textbf{0.00730} & $\sim 0\%$ \\
6.25 & 23.75 & 0.007091 & $-2.8\%$ \\
6.5 & 24.70 & 0.007063 & $-3.2\%$ \\
7 & 26.39 & 0.006797 & $-6.9\%$ \\
8 & 30.40 & 0.006400 & $-12.3\%$ \\
9 & 34.20 & 0.006065 & $-16.9\%$ \\
\bottomrule
\end{tabular}
\end{center}

%=============================================================================
\section{Related Numerical Values Appearing Across Papers}
%=============================================================================

\subsection{The Value 3.797 / 3.80}

The shifted Chern--Simons level expectation value at $k_{\max} = 60$:
\begin{equation}\tag{NUM-1}
\langle k + 2 \rangle_{k_{\max}=60} = 3.7969 \approx 3.80.
\end{equation}

This appears in:
\begin{itemize}
\item section\_quantum.tex (Bridge Lemma box)
\item Ab Initio paper (Eqs.~15--16, Table I, all results sections)
\item section\_alpha\_locked.tex (closure discussion)
\end{itemize}

\subsection{The Value 7.25 / 7.27}

The measured $U(1)$ stiffness $\kappa_{U(1)} \approx 7.25$ (Ab Initio Eq.~11).

Note: the value $7.27 \times 10^{-5}$ rad/s appearing in some research output files
is Earth's angular rotation velocity $\omega_\oplus$, used in clock-drift calculations ---
\textbf{not} related to $\kappa_{U(1)}$.

\subsection{The Stiffness Values}

From the lattice at $(\beta_{U(1)}, \beta_{SU(2)}) = (3.80, 22.80)$:
\begin{align}
\kappa_{U(1)} &\approx 7.25, \tag{STIFF-1} \\
\kappa_{SU(2)} &\approx 14.5, \tag{STIFF-2} \\
\kappa_{U(1)}/\kappa_{SU(2)} &= 0.495 \pm 0.020 \approx 1/2. \tag{STIFF-3}
\end{align}

%=============================================================================
\section{Summary: The Full Equation Map}
%=============================================================================

\begin{tcolorbox}[colback=blue!5!white, colframe=blue!50!black, title=Complete $\alpha$ Equation Inventory]

\textbf{Topological inputs (zero free parameters):}
\begin{enumerate}
\item $k_{\max} = \chi(\mathbb{C}P^2, \mathcal{O}(9) \oplus \mathcal{O}^{\oplus 5}) = 60$ \quad (Bridge Lemma)
\item $\beta_{U(1)} = \langle k+2\rangle_{k_{\max}=60} = 3.80$ \quad (CS vacuum)
\item $\beta_{SU(2)}/\beta_{U(1)} = (n_2/n_1) \times N_{\mathrm{gen}} = 2 \times 3 = 6$ \quad (geometry + index)
\item $\kappa_{U(1)}/\kappa_{SU(2)} = n_1/n_2 = 1/2$ \quad (Frame Stiffness Theorem)
\end{enumerate}

\textbf{Extraction formula:}
\begin{equation*}
\alpha_W = \frac{(1/\kappa_{U(1)})(4/\kappa_{SU(2)})}{(1/\kappa_{U(1)}) + (4/\kappa_{SU(2)})} \cdot \frac{1}{4\pi}
\end{equation*}

\textbf{Convention-locked precision:}
\begin{equation*}
\alpha^{-1} = 137.03599985 \quad (-0.006\;\text{ppm})
\end{equation*}

\textbf{Physical consequences:}
\begin{itemize}
\item $k_a = 3/(8\alpha) \approx 51.4$ \quad (self-coupling)
\item $k_\alpha = \alpha^2/(2\pi)$ \quad (clock coupling)
\item $a_0 = 2\sqrt{\alpha}\,cH_0$ \quad (MOND scale)
\item $v = M_P \alpha^8 \sqrt{2\pi}$ \quad (electroweak hierarchy)
\item $m_f = A_f \alpha^{n_f} v/\sqrt{2}$ \quad (fermion masses)
\end{itemize}
\end{tcolorbox}

\end{document}
